%% ============================================================
%%  Вопрос 2: k-связные графы; граф блоков-точек сочленения
%% ============================================================
\section{\texorpdfstring{$k$}{k}-связные графы; граф блоков-точек сочленения}

\subsection{\texorpdfstring{$k$}{k}-связные графы}

\begin{Definition}{$k$-связный граф}{}
  Граф $G$ называется \emph{$k$-связным}, если
  $|V(G)| > k$ и удаление любого множества $S \subset V(G)$
  с $|S| < k$ оставляет граф связным.
\end{Definition}

\begin{Definition}{}{}
  $\kappa_G(v, u)$ — максимальное количество попарно непересекающихся
  по внутренним вершинам путей от $v$ до $u$ в графе $G$.
\end{Definition}

\begin{Theorem}{}{}
  Граф $G$ является $k$-связным $\Leftrightarrow$
  $\kappa_G(v, u) \geq k$ \; $\forall\, v \neq u$.
\end{Theorem}

\begin{proof}
  \subsubsection*{$(\Rightarrow)$}

  Пусть $G$ $k$-связен; возьмём произвольные $v \neq u$.

  \begin{itemize}
    \item \textit{$v$ и $u$ несмежны.}
          $\kappa_G(v,u) \geq k$ непосредственно из теоремы Менгера.

    \item \textit{$v$ и $u$ смежны.}
          Добавим вершину $v'$, соединённую с соседями $u$,
          и $u'$, соединённую с соседями $v$.
          Связность не понизится, поэтому по теореме Менгера
          существует $k$ непересекающихся путей от $v'$ до $u'$.
          Каждый такой путь заменяется на путь $v\!\to\!u$:
          \begin{enumerate}
            \item содержит $v$ и $u$ $\Rightarrow$ ребро $\{v,u\}$;
            \item содержит только $v$ $\Rightarrow$ заменяем $u'$ на $u$;
            \item содержит только $u$ $\Rightarrow$ заменяем $v'$ на $v$;
            \item не содержит ни $v$, ни $u$ $\Rightarrow$ заменяем $v'\!\to\!v$, $u'\!\to\!u$.
          \end{enumerate}
  \end{itemize}

  \subsubsection*{$(\Leftarrow)$}

  Из $k$ внутренне непересекающихся путей (не более одного из которых
  является ребром $\{v,u\}$, а остальные содержат хотя бы одну
  внутреннюю вершину) следует $|V(G)| \geq k+1 > k$.

  Предположим, $\exists\, S$ с $|S| < k$ такое, что $G\setminus S$ несвязен.
  Тогда $\exists\, a,b \in V\setminus S$ в разных компонентах,
  то есть $\kappa_G(a,b) \leq |S| < k$ — противоречие.
\end{proof}

\begin{Example}{Граф $C_5$ — 2-связный, но не 3-связный}{}
  Цикл $C_5$ на $\{1,2,3,4,5\}$:
  \begin{itemize}
    \item Удаление любой одной вершины оставляет путь $P_4$ — связный.
    \item Для любых $u\neq v$: два пути по дугам цикла внутренне не пересекаются,
          поэтому $\kappa_{C_5}(u,v)=2$.
    \item Удаление $\{1,3\}$ отделяет вершину 2, значит $\kappa(C_5)=2$
          и $C_5$ не является $3$-связным.
  \end{itemize}
\end{Example}

\subsection{Точки сочленения и блоки}

\begin{Definition}{}{}
  Вершина $v \in V(G)$ называется \emph{точкой сочленения},
  если граф $G \setminus \{v\}$ несвязен.
\end{Definition}

\begin{Definition}{}{}
  Индуцированный подграф $B$ называется \emph{блоком}, если:
  \begin{itemize}
    \item $B$ — ребро или двусвязный граф;
    \item $B$ максимален по включению среди таких подграфов.
  \end{itemize}
\end{Definition}

\begin{Statement}{}{}
  Если два различных блока $B_1$ и $B_2$ пересекаются,
  то $|B_1 \cap B_2| = 1$, и единственная общая вершина
  является точкой сочленения.
\end{Statement}

\begin{proof}
  Из максимальности $B_1 \cup B_2$ не является блоком,
  поэтому $|B_1 \cap B_2| \le 1$.
  Пусть $B_1 \cap B_2 = \{v\}$.
  Если $v$ — не точка сочленения, то между любыми вершинами
  $B_1 \cup B_2$ есть путь, избегающий $v$, и $B_1 \cup B_2$
  двусвязен — противоречие максимальности.
\end{proof}

\begin{figure}[H]
\centering
\begin{tikzpicture}[
  >=Stealth, font=\small,
  v/.style={circle, draw, fill=white, minimum size=8mm, inner sep=0pt},
  vc/.style={circle, draw=orange!80!black, fill=orange!25, minimum size=8mm,
             inner sep=0pt, line width=1.5pt}]

\node[v]  (v1) at (-1.5,  1) {$1$};
\node[v]  (v2) at (-1.5, -1) {$2$};
\node[vc] (v3) at ( 0,    0) {$3$};
\node[v]  (v4) at ( 1.8,  1) {$4$};
\node[vc] (v5) at ( 1.8, -1) {$5$};
\node[v]  (v6) at ( 3.6,  0) {$6$};

\draw (v1)--(v2) (v1)--(v3) (v2)--(v3);
\draw (v3)--(v4) (v4)--(v5) (v5)--(v3);
\draw (v5)--(v6);

\begin{scope}[on background layer]
  \fill[blue!12, rounded corners=10pt]
    (-2.1,-1.6) -- (-2.1,1.6) -- (0.7,1.6) -- (0.7,-1.6) -- cycle;
  \fill[green!12, rounded corners=10pt]
    (-0.6,-1.6) -- (-0.6,1.6) -- (2.5,1.6) -- (2.5,-1.6) -- cycle;
  \fill[violet!12, rounded corners=8pt]
    (1.1,-0.6) -- (1.1,0.6) -- (4.2,0.6) -- (4.2,-0.6) -- cycle;
\end{scope}

\node[blue!70!black]   at (-1.3, 2.1) {$B_1$};
\node[green!60!black]  at ( 1.0, 2.1) {$B_2$};
\node[violet!60!black] at ( 3.2, 1.1) {$B_3$};
\node[orange!80!black, font=\footnotesize] at (1.0, -2.1)
  {оранжевые вершины — точки сочленения};
\end{tikzpicture}
\caption{Граф с блоками $B_1=\{1,2,3\}$, $B_2=\{3,4,5\}$, $B_3=\{5,6\}$
  и точками сочленения $3$ и $5$.
  Цветные области перекрываются в точках сочленения — они входят в оба соседних блока.}
\end{figure}

\subsection{Дерево блоков–точек сочленения}

Обозначим $A$ — множество точек сочленения, $\mathcal{B}$ — множество блоков.
\emph{Дерево блоков и точек сочленения} графа $G$ — граф с вершинами
$A \cup \mathcal{B}$, в котором $a \in A$ и $B \in \mathcal{B}$ соединены ребром
тогда и только тогда, когда $a \in V(B)$.

\begin{Theorem}{}{}
  Описанный граф является деревом.
\end{Theorem}

\begin{proof}
  \begin{itemize}
    \item У двух различных блоков не более одной общей вершины:
          иначе их объединение было бы блоком — противоречие максимальности.
    \item По той же причине в графе нет цикла из блоков.
  \end{itemize}
\end{proof}

\begin{figure}[H]
\centering
\begin{tikzpicture}[
  >=Stealth, font=\small,
  vcut/.style={circle, draw=orange!80!black, fill=orange!25, minimum size=10mm,
               inner sep=0pt, line width=1.5pt},
  blk/.style={rectangle, draw, rounded corners=5pt, minimum width=14mm,
              minimum height=10mm, inner sep=3pt, font=\bfseries}]

\node[blk, draw=blue!60!black,   fill=blue!12]   (B1) at (0,   0) {$B_1$};
\node[vcut]                       (c3) at (2.2, 0) {$3$};
\node[blk, draw=green!60!black,  fill=green!12]  (B2) at (4.4, 0) {$B_2$};
\node[vcut]                       (c5) at (6.6, 0) {$5$};
\node[blk, draw=violet!60!black, fill=violet!12] (B3) at (8.8, 0) {$B_3$};

\draw (B1)--(c3)--(B2)--(c5)--(B3);
\end{tikzpicture}
\caption{Дерево блоков–точек сочленения для графа выше.
  Прямоугольники — блоки, кружки — точки сочленения.
  Граф связен и ацикличен — дерево.}
\end{figure}
